\subsection{Asymmetric Observations}
\label{app:rl_training_details}

Table~\ref{tab:teacher_observations} contrasts the actor's deployable noisy
estimates with the centralized critic's privileged simulator state.  Running
mean/variance normalization is applied to both inputs and to value targets.

\begin{table}[H]
  \centering
  \caption{Asymmetric teacher observations.}
  \label{tab:teacher_observations}
  \footnotesize
  \setlength{\tabcolsep}{3pt}
  \renewcommand{\arraystretch}{1.0}
  \begin{tabularx}{\columnwidth}{@{}l r X@{}}
    \toprule
    Input & Dim. & Component \\
    \midrule
    \multirow{3}{*}{\makecell[l]{Actor\\(87)}}
      & 26 & Noisy actuated-joint positions and velocities \\
      & 36 & Noisy palm/fingertip positions and velocities \\
      & 25 & Noisy object pose (7), goal position (3), object ID and scale (2), and previous action (13) \\
    \midrule
    \multirow{4}{*}{\makecell[l]{Critic\\(133)}}
      & 26 & Clean actuated-joint positions and velocities \\
      & 54 & Palm/fingertip positions and spatial velocities \\
      & 22 & Palm contact force and measured joint torques \\
      & 31 & Object pose and spatial velocity (13), goal position (3), object ID and scale (2), and previous action (13) \\
    \bottomrule
  \end{tabularx}
\end{table}

\subsection{VLM Prompt}
\label{app:vlm_inference_details}

A typical prompt specifies the system and user requirements, intervention rule,
multimodal input format, task description, and output keys, as shown in the
example below.

\begin{algorithm}[H]
\caption*{\textbf{Typical VLM Prompt Example}}
{\footnotesize\ttfamily
\begin{algorithmic}
\STATE \textbf{System:}
\STATE You are a robotics safety assistant. Your task is to assess the current
safety--autonomy operating point of a visuomotor policy.
\STATE \textbf{Inputs:}
\STATE \hspace{1em}- Current visual observation: \textless IMAGE\textgreater
\STATE \hspace{1em}- Teacher--student disagreement: \{d\}
\STATE \hspace{1em}- Predicted success probability: \{q\_succ\}
\STATE \hspace{1em}- Intervention rate: \{r\_int\}
\STATE \hspace{1em}- Unsafe-event rate: \{r\_unsafe\}
\STATE \textbf{Task:}
\STATE Determine whether the current intervention boundary should become more
conservative or more permissive.
\STATE \textbf{Return:}
\STATE \{
\STATE \hspace{1em}"disagreement\_threshold": \ldots,
\STATE \hspace{1em}"success\_threshold": \ldots,
\STATE \hspace{1em}"reason": \ldots\ \}
\end{algorithmic}
}
\end{algorithm}

\subsection{Additional Experimental Results}
\label{app:additional_experimental_results}

We will report the additional real-world evaluation using the same success and
safety breakdown as Table~\ref{tab:real_world_results}. Each column represents
one test object, while the rows separate task success from the observed unsafe
failure modes. The object names and numerical results are retained as
placeholders pending the expanded evaluation.

\begin{table}[H]
\centering
\small
\caption{Template for additional experimental results.  Object names and
numerical entries are placeholders.}
\label{tab:additional_experimental_results}
\setlength{\tabcolsep}{4pt}
\renewcommand{\arraystretch}{1.2}
\resizebox{0.48\textwidth}{!}{%
\begin{tabular}{l|cccccc}
\hline
Metric (\%) & Object 1 & Object 2 & Object 3 & Object 4 & Object 5 & Object 6 \\
\hline
Lift Success      & -- & -- & -- & -- & -- & -- \\
Unsafe Rate Avg   & -- & -- & -- & -- & -- & -- \\
Harmful collision & -- & -- & -- & -- & -- & -- \\
Object escape     & -- & -- & -- & -- & -- & -- \\
Palm flipped      & -- & -- & -- & -- & -- & -- \\
Hand too far      & -- & -- & -- & -- & -- & -- \\
\hline
\end{tabular}%
}
\end{table}
